\section{Introduction}
\label{sec:intro}

Stablecoins are becoming payment infrastructure, and the GENIUS Act gives them their
first comprehensive U.S. federal framework~\cite{geniusact2025}. A fiat-backed stablecoin that
trades below par admits two readings. Under rational repricing, an impaired reserve
is worth less. Under non-fundamental repricing, a self-fulfilling coordination
run~\cite{diamonddybvig1983,goldsteinpauzner2005}
and microstructure friction push the price below the floor a reserve loss can justify. No
measurement of March 2023 cited in Section~\ref{sec:lit} forms that floor. A de-peg with
a non-fundamental share bounded away from zero calls for backstop and redemption design,
since a reserve-quality rule addresses only the fundamental share.

We offer a reusable estimator for any disclosed-impairment de-peg. From an issuer's
disclosed reserve impairment, derive a worst-case floor. Test whether the observed
trough breaches that floor. Read the breach, if any, as a bound on the non-fundamental
share. Figure~\ref{fig:schematic} shows those three steps. A supervisor can run it on any
public hourly print, within hours of a break. It returns no verdict when the trough does not
breach the floor.

%% H2 T6. Declared inside Section 1, immediately after the paragraph that cites it,
%% rather than in the skeleton between \input lines. A float is only considered from
%% its own position forward, and this is the earliest position that is still after its
%% first citation. The fragment is \input, not \includegraphics: it is a bare
%% tikzpicture with no figure environment of its own, so the wrapper lives here.
\begin{figure}[t]
\centering
\input{figures/fig_schematic}
\caption{The estimator. A disclosed impairment gives a floor, and a public price that
breaches it bounds the non-fundamental share.}
\label{fig:schematic}
\end{figure}

We apply it to the March 2023 USD~Coin (USDC) de-peg, where Silicon Valley Bank (SVB)
failed holding about $8\%$ of reserves, and a joint federal
backstop~\cite{fedtreasfdic2023pressrelease} restored the peg within days.

Under expected-loss pricing, fundamental value cannot fall below $1-\phi$, where $\phi$ is
the disclosed impairment fraction. Under (A4), a preference-free corollary extends that
floor to any monotone-utility valuation. The observed trough breaches the floor at every admissible
unremediated-loss probability and recovery rate, and requires $q^{\star}=1.54$, which
exceeds one. Under (A1)--(A3), no expected-loss account of the disclosed impairment is
consistent with the trough. At the disclosed $\phi$ the lower edge of the resulting
non-fundamental share runs from $19\%$ on Kraken's close to $35\%$ on the composite. It
stays positive across impairment readings, and at the full re-peg is $100\%$ ex post, an
identity
(Section~\ref{sec:ident}).

The paper makes two contributions:
\begin{itemize}
\item \textbf{An identification argument} (Proposition~\ref{prop:floor}) turning a
  disclosed impairment into a floor that a price either breaches or does not, with the
  breach read as a bound on the non-fundamental share.
\item \textbf{A screen} over six candidate de-pegs, none of which satisfies all three
  preconditions, so the single qualifying application is a property of the disclosure
  record rather than of the method.
\end{itemize}
Two analyses stay outside the estimator: the ex-ante fragility
ordering by reserve exposure that the realized troughs match, with null probability $1/2$
(Section~\ref{sec:placebo}), and a three-episode panel (Section~\ref{sec:spec}).
